\documentclass[12pt,a4paper]{article}
\usepackage{goyabase}

\usepackage{longtable}
\usepackage{calc}

% Pandoc compatibility
\providecommand{\tightlist}{%
  \setlength{\itemsep}{0pt}\setlength{\parskip}{0pt}}
\newcommand{\passthrough}[1]{#1}

\title{\textbf{Más ejemplos Redis - GloboTour}}
\author{Víctor de Juan}
\date{\today}

\usepackage[spanish, es-noshorthands, es-tabla]{babel}
\usepackage[autostyle]{csquotes}
\begin{document}

\maketitle

\hypertarget{solucionario-tarea-1-redis---globotour}{%
\section{Solucionario Tarea 1: Redis -
GloboTour}\label{solucionario-tarea-1-redis---globotour}}

Este documento contiene los comandos necesarios para resolver la Tarea 1
de la práctica NoSQL.

\hypertarget{gestiuxf3n-de-sesiones-strings}{%
\subsection{Gestión de Sesiones
(Strings)}\label{gestiuxf3n-de-sesiones-strings}}

\textbf{Comando:} Crear un token con expiración de 30 minutos (1800
segundos).

\begin{lstlisting}
SET sesion:token:9988 "user_id:105" EX 1800
\end{lstlisting}

\emph{Para verificar el tiempo restante:}
\texttt{TTL sesion:token:9988}

\hypertarget{muxe9tricas-en-vivo-counters}{%
\subsection{Métricas en Vivo
(Counters)}\label{muxe9tricas-en-vivo-counters}}

\textbf{Comando:} Inicializar e incrementar visitas.

\begin{lstlisting}
# Incrementamos en 3 directamente o 3 veces con INCR
INCRBY metricas:visitas:oferta:paris 3
\end{lstlisting}

\hypertarget{ranking-de-popularidad-sorted-sets}{%
\subsection{Ranking de Popularidad (Sorted
Sets)}\label{ranking-de-popularidad-sorted-sets}}

\textbf{Comando:} Crear el top de ventas y recuperar los 2 mejores.

\begin{lstlisting}
ZADD experiencias:top_ventas 150 "Vuelo en Globo Segovia" 320 "Cata de Vinos Ribera" 95 "Buceo Cabo de Gata" 210 "Escape Room Medieval"

# Recuperar Top 2 (descendente) con puntuación
ZREVRANGE experiencias:top_ventas 0 1 WITHSCORES
\end{lstlisting}

\hypertarget{categorizaciuxf3n-sets}{%
\subsection{Categorización (Sets)}\label{categorizaciuxf3n-sets}}

\textbf{Comando:} Crear conjuntos de etiquetas y buscar la común.

\begin{lstlisting}
SADD user:105:intereses "aventura" "aire_libre" "fotografia"
SADD user:202:intereses "gastronomia" "vinos" "relax" "aventura"

# Intersección para ver qué tienen en común
SINTER user:105:intereses user:202:intereses
\end{lstlisting}

\hypertarget{carga-y-anuxe1lisis-del-catuxe1logo-hashes}{%
\subsection{Carga y Análisis del Catálogo
(Hashes)}\label{carga-y-anuxe1lisis-del-catuxe1logo-hashes}}

\textbf{Recordando: comando de carga (desde la terminal de Linux):}

\begin{lstlisting}[language=bash]
# Estando en la carpeta redis/
docker exec -i redis-server redis-cli < src/datos_iniciales.redis
\end{lstlisting}
\textbf{Consultar la experiencia 500:}

\begin{lstlisting}
HGETALL exp:500
\end{lstlisting}

\textbf{Crear Ranking ``Premium'' basado en precios:} Para este paso, primero consultamos los precios reales de las experiencias en el sistema:

\begin{lstlisting}
HGET exp:500 precio  -- Devuelve 150
HGET exp:501 precio  -- Devuelve 45
HGET exp:502 precio  -- Devuelve 80
HGET exp:503 precio  -- Devuelve 25
\end{lstlisting}

Una vez conocidos los precios, creamos el Sorted Set ordenado por dicho valor:

\begin{lstlisting}
ZADD catalogo:premium 150 "exp:500" 45 "exp:501" 80 "exp:502" 25 "exp:503"

# Visualizar de más caro a más barato
ZREVRANGE catalogo:premium 0 -1 WITHSCORES
\end{lstlisting}

\end{document}
